% Suggested final budget: 3.1 pages, including plots and tables:
% 4.1 = 1.2, 4.2 = 0.9, 4.3 = 0.6, 4.4 = 0.4.
\section{Results}
\label{sec:results}
\subsection{Six-task baselines and effect of demonstration count}
\label{sec:datasize}
% Values transcribed 2026-10-12 from dp-manip/docs/results-summary.zh-CN.md
% (updated 2026-10-07; sections 2 and 7.4), a preliminary summary of the member
% reports. PushCube means and SDs, and the success_at_end SDs, are computed from
% the per-seed values recorded in dp-manip/scripts/plot_cross_task.py.
% Verify every cell against the raw evaluation output before submission.
% The PlaceSphere row was checked on 2026-10-13 against the raw evaluation files
% (~/dp-runs-placesphere/*/eval/test_final_h200.json on the cluster): all 21
% per-run success counts match, and success_at_end equals success_once in
% every run.
% Still open: PushCube success_at_end SD; LiftPegUpright success_at_end under
% the 200-step horizon used in this table.
% Keep in sync with scripts/plot_datasize.py.
\begin{table}[htbp]
  \centering
  \caption{Six-task UNet baselines across demonstration counts $N$:
  \metriconce{} (\%) of the final checkpoint on 100 held-out test episodes per
  training seed, mean $\pm$ sample standard deviation over training seeds
  (3 for $N\le 50$, 5 for $N\ge 100$). The last column (End) is
  \metricend{} at $N=100$. ``--'': not run, because the pre-registered rule
  adds $N=400$ only when the $100\rightarrow200$ gain is distinguishable from
  zero. \textit{TBD}: not yet filled.}
  \label{tab:datasize}
  \small
  \setlength{\tabcolsep}{4pt}
  \begin{tabularx}{\linewidth}{@{}Ycccccc@{}}
    \toprule
    & \multicolumn{5}{c}{\metriconce{} (\%) at $N$} & End (\%) \\
    \cmidrule(lr){2-6}\cmidrule(l){7-7}
    Task & 25 & 50 & 100 & 200 & 400 & $N=100$ \\
    \midrule
    PickCube & $0.7\pm1.2$ & $0.7\pm0.6$ & $19.2\pm1.8$ & $45.2\pm3.1$ & $53.4\pm2.1$ & $12.8\pm0.8$ \\
    StackCube & $1.7\pm1.2$ & $12.3\pm2.5$ & $41.0\pm2.9$ & $66.0\pm3.1$ & $83.2\pm1.9$ & $35.6\pm2.5$ \\
    PushCube & $45.3\pm3.2$ & $82.3\pm1.5$ & $94.8\pm1.3$ & $96.0\pm1.9$ & -- & 54$^{a}$ \\
    PullCube & $76.7\pm1.2$ & $91.3\pm0.6$ & $95.6\pm0.9$ & $97.2\pm0.4$ & -- & $95.6\pm0.9$ \\
    PlaceSphere & $4.7\pm2.9$ & $25.3\pm4.1$ & $36.0\pm3.6$ & $76.6\pm5.6$ & $96.2\pm0.8$ & $36.0\pm3.6$ \\
    LiftPegUpright$^{b}$ & $7.0\pm1.0$ & $36.3\pm3.1$ & $67.4\pm2.1$ & $74.8\pm2.8$ & $79.2\pm0.8$ & \pending \\
    \bottomrule
  \end{tabularx}

  \vspace{2pt}
  \parbox{\linewidth}{\footnotesize
  $^{a}$ Mean only; seed-level values not yet reported.\quad
  $^{b}$ Evaluated with a 200-step horizon; \cref{tab:backbones} uses 400.}
\end{table}

% Sections 4.1 and 4.2 drafted 2026-10-12 from
% dp-manip/docs/results-summary.zh-CN.md (sections 2, 3, 7) and
% report/report_pickcube_yhr.md (PickCube intervals, validation loss, cost).
% UNVERIFIED: the numbers in tab:datasize, tab:backbones, tab:takeover,
% tab:moore and in the text of this section are transcribed from the group's
% results summary and member reports. Check them against raw evaluation output
% before submission.

\textbf{Baselines.} With the shared configuration and $N=100$, the UNet
baseline solves the two non-prehensile tasks in about 95\% of test episodes,
LiftPegUpright in two thirds, StackCube and PlaceSphere in roughly 40\%, and
PickCube in 19\% (\cref{tab:datasize}). \metricend{} matches \metriconce{} on
PullCube but is far lower on PushCube (54\% against 95\%): the cube reaches
the goal region and then slides out.

% Generated by scripts/plot_datasize.py from the same numbers as tab:datasize.
\begin{figure}[htbp]
  \centering
  \includegraphics[width=\linewidth]{datasize_curves.pdf}
  \caption{\metriconce{} of the UNet baseline on held-out test seeds against
  the number of training demonstrations. Points are means over training seeds;
  error bars are $\pm1$ sample standard deviation. PushCube and PullCube stop at
  $N=200$ under the pre-registered rule. PegInsertionSide is measured on 50
  validation episodes. Control modes and horizons differ between tasks
  (\cref{tab:tasks}), so curve shapes are comparable but levels are not.}
  \label{fig:data-efficiency}
\end{figure}
\textbf{The data requirement depends on the task.} \Cref{fig:data-efficiency}
shows three patterns. PushCube and PullCube are saturated by $N=100$: doubling
to 200 changes success by $+1.2$ points [$-1.4$, $+3.8$] and $+1.6$
[$-0.2$, $+3.6$], and PullCube already reaches 76.7\% with 25 demonstrations.
StackCube and PlaceSphere are still climbing steeply at $N=200$. StackCube
gains $+25.0$ [$+18.2$, $+31.8$] from 100 to 200 and $+17.2$
[$+10.8$, $+23.6$] from 200 to 400, ending at 83.2\%; PlaceSphere gains
$+40.6$ [$+32.6$, $+48.2$] and then $+19.6$ [$+13.6$, $+25.8$], ending at
96.2\%. PickCube and LiftPegUpright show diminishing
returns below the ceiling: PickCube gains $+26.0$ [$+19.8$, $+32.2$] and then
only $+8.2$ [$+1.4$, $+15.0$], stopping at 53.4\%. On the logarithmic axis the
curves have a similar shape and differ mainly in position: by interpolation,
50\% success needs fewer than 25 demonstrations on PullCube, about 27 on
PushCube, 68 on LiftPegUpright, and close to 130 on StackCube and PlaceSphere.

Two qualifications apply. PickCube needs about 300 demonstrations by the same
measure, more than StackCube, which contains a pick; \cref{sec:backbone} shows
that this reflects the UNet on this task, so the curves describe this
configuration, not the tasks in general. And PegInsertionSide stays near zero
at every size.
% The PlaceSphere intervals above were computed for this report with
% scripts/bootstrap_adjacent.py from the per-seed rates in
% dp-manip/docs/experiment-progress.zh-CN.md sections 6 and 15. The same script
% reproduces the PickCube intervals of the member report.
\todo{Add the adjacent-$N$ intervals for LiftPegUpright once its per-seed
results are available; the member report gives only means.}

\subsection{Backbone success and efficiency}
\label{sec:backbone}
\begin{table}[htbp]
  \centering
  \caption{Placeholder for a backbone comparison on one task and one dataset
  size. Specify task, $N$, training seed count, evaluation episodes, GPU model,
  and whether inference timing measures a denoising step or a full action chunk.
  Repeat task blocks as needed; do not average incompatible hardware timings.}
  \label{tab:backbones}
  \small
  \begin{tabularx}{\linewidth}{@{}Yccccc@{}}
    \toprule
    Backbone & \shortstack{Once\\(\%)} & \shortstack{Params\\(M)} &
    \shortstack{Train\\(ms/step)} & \shortstack{VRAM\\(GB)} &
    \shortstack{Infer\\(ms)} \\
    \midrule
    UNet & \pending & \pending & \pending & \pending & \pending \\
    Transformer & \pending & \pending & \pending & \pending & \pending \\
    MLP & \pending & \pending & \pending & \pending & \pending \\
    \bottomrule
  \end{tabularx}
\end{table}

\textbf{No backbone is best on every task.} The Transformer is best on
PickCube and PullCube ($+2.4$ points over the UNet on PullCube,
distinguishable) and worst on PushCube ($-5.8$, distinguishable). The UNet is
best or tied on StackCube, PushCube, and LiftPegUpright. The MLP is within
about ten points of the best backbone on the three 100-step tasks but clearly
worst on StackCube (25.0\% against 41.0\%) and LiftPegUpright
($-30.4$ points [$-40.2$, $-21.2$] against the UNet). The Transformer is the
most consistent: it stays within 90\% of the best backbone on every task, and
its deficit on LiftPegUpright ($-6.8$ [$-14.2$, $+0.8$]) is not
distinguishable from zero.

\textbf{PickCube is the outlier.} The UNet reaches 19.2\%, against 90.6\% for
the Transformer ($+71.4$ points [$+65.6$, $+76.8$]) and 79.6\% for the MLP
($+60.4$ [$+55.0$, $+65.6$]). Its validation denoising loss rises from 0.018
at 10k steps to 0.064 at 100k while its training loss falls to
$5\times10^{-4}$; the other two backbones end at 0.028--0.030. The final UNet
checkpoint is therefore overfitted on this task, and changing the backbone at
$N=100$ helps more than quadrupling the data for the UNet (53.4\% at $N=400$).
We do not know why this happens only on PickCube: the UNet also drives the
training loss to near zero on StackCube and LiftPegUpright without the same
collapse.

\textbf{Holding the success state.} Ranking by \metricend{} changes the order
on two tasks. On PushCube the MLP keeps the cube in the goal region most often
(67\%, against 63\% for the Transformer and 54\% for the UNet), although the
UNet leads on \metriconce{}. On LiftPegUpright the Transformer leads
($55.6\pm5.4$\% against $49.2\pm12.7$\% for the UNet), and the MLP raises the
peg in 41.0\% of episodes but keeps it upright in only 5.2\%. The Transformer
is never the worst backbone by this metric.

\textbf{Cost.} On PickCube with one RTX~4080 SUPER, a 100k-step run takes
1.72~h for the UNet, 1.75~h for the Transformer, and 0.74~h for the MLP; one
policy call of 100 denoising steps takes 540, 570, and 97~ms; and the
denoisers have 66.4M, 8.97M, and 0.353M parameters. The MLP's deficit appears
on the two tasks that are both harder and longer, so these two explanations
cannot be separated here; further limits are listed in \cref{sec:limits}.

\subsection{Expert takeover correction}
\label{sec:takeover}
% Drafted 2026-10-12. Sources: dp-manip/docs/failure-expert-takeover-probe.zh-CN.md
% sections 7 (first probe, failure stages), 10 (main result), 12 (dose and
% mechanism), and dp-manip/docs/results-summary.zh-CN.md section 4 (guidance).
% Verify against results/takeover-correction/decision.json and
% results/takeover-curve/curve_decision.json.
% Values transcribed 2026-10-12 from section 4 of
% dp-manip/docs/results-summary.zh-CN.md (updated 2026-10-07). Checked on
% 2026-10-13 against results/takeover-correction/decision.json: the per-arm
% counts and the three differences with their intervals match.
\begin{table}[htbp]
  \centering
  \caption{Takeover-correction study on PlaceSphere, starting from the three
  $N=100$ UNet baseline checkpoints: 200 fresh evaluation episodes per
  checkpoint and arm (600 per arm). $\Delta$ is the paired difference to arm B
  in percentage points with its 95\% bootstrap interval. C$-$D is
  $+14.3$ [$+9.5$, $+19.0$].}
  \label{tab:takeover}
  \small
  \begin{tabularx}{\linewidth}{@{}lYccc@{}}
    \toprule
    Arm & Fine-tuning data & \shortstack{Successes /\\episodes} &
    \shortstack{\metriconce{}\\(\%)} & \shortstack{$\Delta$ vs.\ B\\(pp) [95\% CI]} \\
    \midrule
    B & None (baseline checkpoint) & 197 / 600 & 32.8 & -- \\
    C & Base demos + takeover corrections & 280 / 600 & 46.7 & $+13.8$ [$+9.3$, $+18.3$] \\
    D & Base demos + matched extra demos & 194 / 600 & 32.3 & $-0.5$ [$-4.7$, $+3.5$] \\
    \bottomrule
  \end{tabularx}
\end{table}

\textbf{Where the baseline fails.} On PlaceSphere at $N=100$ the baseline
fails in about two thirds of episodes, almost always at the grasp. Of 338
failures in 488 rollouts of one checkpoint, 49\% never grasped the sphere and
41\% touched it without holding a stable grasp. A first design that took over after a stable
grasp therefore applied to only about 9\% of failures and was stopped. Moving the takeover to the missed grasp
covered 16 of 18 baseline failures in a 30-seed probe, and the scripted expert
recovered 12 of those 16.

\textbf{Corrections versus more demonstrations.} Fine-tuning on 100
corrections raises success from 32.8\% to 46.7\% (\cref{tab:takeover}), in the
same direction for all three checkpoints, while a matched amount of ordinary
demonstrations leaves it unchanged. The gain grows with the number of
corrections ($+4.3$ points [$+0.7$, $+7.8$] with 25, $+6.0$ [$+1.7$, $+10.8$]
with 50), whereas ordinary demonstrations help only in much larger amounts
($+18.2$ [$+6.2$, $+33.2$] with 300 extra), roughly four times the training
windows for the same gain. The two kinds of data also act differently: after
a missed grasp the corrections arm still succeeds in 10.7\% of episodes,
against at most 0.3\% for every other arm, whereas 300 extra demonstrations
make the policy miss less often (63.7\% to 46.7\% of episodes).

\textbf{Scope.} An earlier attempt that used failed rollouts to steer sampling
away from failure did not help ($+1.7$ points [$-2.0$, $+5.0$]) and doubled
inference time. Corrections are not a substitute for data at scale: training
from scratch on 200 demonstrations reaches 76.6\% (\cref{tab:datasize}). The
study covers one task, one baseline level, and three checkpoints; the trigger
and the expert read privileged simulator state; and recovery after a missed
grasp remains rare.

\subsection{Moore Threads profiling and benchmarking}
\label{sec:moore}
% Drafted 2026-10-12 from report/moer_pickcube_backbone_runbook.md on branch
% origin/moore-transformer-attn-fix (sections 2, 6, 7, 8) and commit 693ec7f.
% That branch is not merged here; verify the numbers against the summary.json
% and eval files on the Moore Threads server before submission.
% Sources:
% - RTX 4080 SUPER training time: report/report_pickcube_yhr.md section 5.
% - RTX 4080 success, seeds 1-3: per-seed values in report/report_pickcube_yhr.md
%   sections 3.2 and 4 (UNet 0.17/0.18/0.19, Transformer 0.95/0.92/0.92,
%   MLP 0.81/0.79/0.82).
% - MTT S4000: report/moer_pickcube_backbone_runbook.md on branch
%   origin/moore-transformer-attn-fix, sections 8.1, 8.2 and 8.5. The UNet and
%   MLP training times are read from the completion timestamps in section 8.1
%   (about 4.3 h and 2.0 h per run); the Transformer's 3.7 h is reported
%   directly for the retrained runs. Replace them with summary.json wall times.
\begin{table}[htbp]
  \centering
  \caption{PickCube at $N=100$ on an RTX~4080 SUPER and on a Moore Threads
  MTT~S4000: wall-clock time of one 100k-step training run and \metriconce{}
  (\%) on the 100 test episodes, mean $\pm$ sample standard deviation over
  training seeds 1--3. The S4000 runs train in full precision and the RTX runs
  in mixed precision.}
  \label{tab:moore}
  \small
  \begin{tabularx}{\linewidth}{@{}Ycccc@{}}
    \toprule
    & \multicolumn{2}{c}{Train 100k steps (h)} & \multicolumn{2}{c}{\metriconce{} (\%)} \\
    \cmidrule(lr){2-3}\cmidrule(l){4-5}
    Backbone & RTX 4080 SUPER & MTT S4000 & RTX 4080 SUPER & MTT S4000 \\
    \midrule
    UNet & 1.72 & 4.3 & $18.0\pm1.0$ & $15.3\pm4.9$ \\
    Transformer & 1.75 & 3.7 & $93.0\pm1.7$ & $90.3\pm0.6$ \\
    MLP & 0.74 & 2.0 & $80.7\pm1.5$ & $80.0\pm1.0$ \\
    \bottomrule
  \end{tabularx}
\end{table}

We repeated the PickCube backbone comparison ($N=100$, training seeds 1--3,
same code and protocol) on a Moore Threads MTT~S4000 (48~GB) with PyTorch~2.9
on the MUSA backend (\cref{tab:moore}). Success rates agree with the RTX~4080
runs of the same seeds to within 3 points; the S4000 runs are lower in 7 of 9
pairs, by 2.0 points on average. A 100k-step run takes 2.1 to
2.7 times longer than on the RTX~4080 SUPER. The timing is not like for like:
our trainer enabled mixed precision only for devices reported as CUDA, so the
S4000 runs trained in full precision. One change to the model code was needed.
With the backend's fused scaled-dot-product attention the Transformer's
training loss stalled near 1.48 and all three seeds scored 0 of 100; an
explicit attention implementation that shares the same parameters, with
outputs equal to within $4\times10^{-7}$, restored training, and
\cref{tab:moore} reports the retrained runs.
\FloatBarrier
